%%%%%%%%%%%%%%%%%%%%%%%%%%%%%%%%%%%%%%%%%%%%%%%%%%%%%%%%%%%%%%%%%%%%%%%%%%%%%%%
% IEEE Conference Paper
% IEEEtran class — conference mode
% Author: Anthony Hanna, Johns Hopkins University
%%%%%%%%%%%%%%%%%%%%%%%%%%%%%%%%%%%%%%%%%%%%%%%%%%%%%%%%%%%%%%%%%%%%%%%%%%%%%%%

\documentclass[conference]{IEEEtran}

% Core packages
\usepackage{amsmath,amssymb,amsfonts}
\usepackage{graphicx}
\usepackage{booktabs}
\usepackage{hyperref}
\usepackage{cite}
\usepackage{xcolor}
\usepackage{url}

%%%%%%%%%%%%%%%%%%%%%%%%%%%%%%%%%%%%%%%%%%%%%%%%%%%%%%%%%%%%%%%%%%%%%%%%%%%%%%%
\begin{document}

\title{From Natural Language to Sprint-Ready: A Skill-Based AI Orchestration
Framework for Project Management in GitHub-Native Software Teams}

\author{%
  \IEEEauthorblockN{Anthony Hanna}
  \IEEEauthorblockA{%
    Whiting School of Engineering\\
    Johns Hopkins University\\
    Baltimore, MD, USA\\
    ahanna13@jhu.edu
  }
}

\maketitle

%%%%%%%%%%%%%%%%%%%%%%%%%%%%%%%%%%%%%%%%%%%%%%%%%%%%%%%%%%%%%%%%%%%%%%%%%%%%%%%
\begin{abstract}
Software project management overhead --- ticket creation, triage,
prioritization, and sprint planning --- consumes a disproportionate
share of engineering and product team time, yet remains one of the
least automated phases of the software delivery lifecycle. While
AI-assisted coding tools have demonstrated measurable productivity
gains at the individual contributor level, their application to
structured, workflow-integrated project management remains limited
and largely ad hoc. This paper presents a skill-based AI
orchestration framework that embeds large language model assistance
directly into the full project management lifecycle within
GitHub-native software teams, spanning raw work intake, label
taxonomy-driven triage, capacity-aware sprint planning, workload
rebalancing, and real-time metrics reporting. The framework was
designed and validated through the end-to-end construction of a
full-stack web application, in which the same AI system served as
both development assistant and project management orchestrator.
Structured context files --- encoding team ownership, label ontologies,
sprint capacity, and reusable workflow playbooks --- were used to
ground model outputs and enforce governance checkpoints without
requiring specialized infrastructure or dedicated model deployment.
Compared to ad hoc AI usage, the proposed approach produced more
consistent, enriched, and routing-correct work artifacts while
preserving human oversight at each decision boundary. The results
demonstrate that structured context injection through skill files and
formal label taxonomies is a practical and reproducible pattern for
AI-assisted project management in real-world software teams.
\end{abstract}

\begin{IEEEkeywords}
AI-assisted project management, large language models, GitHub
automation, issue triage, sprint planning, context injection,
skill-based orchestration, software workflow automation
\end{IEEEkeywords}

%%%%%%%%%%%%%%%%%%%%%%%%%%%%%%%%%%%%%%%%%%%%%%%%%%%%%%%%%%%%%%%%%%%%%%%%%%%%%%%
\section{Introduction}
\label{sec:intro}

% TODO: Motivate the problem (PM overhead, gap between coding AI and workflow AI).
% Introduce the framework. State contributions. Outline paper structure.

%%%%%%%%%%%%%%%%%%%%%%%%%%%%%%%%%%%%%%%%%%%%%%%%%%%%%%%%%%%%%%%%%%%%%%%%%%%%%%%
\section{Literature Review}
\label{sec:lit}

% TODO: Survey AI-assisted coding tools (Copilot, Duo, Atlassian Intelligence).
% Cover issue triage research, requirements analysis, RAG, agent orchestration.
% Identify the gap: no integrated, skill-driven PM orchestration layer.

%%%%%%%%%%%%%%%%%%%%%%%%%%%%%%%%%%%%%%%%%%%%%%%%%%%%%%%%%%%%%%%%%%%%%%%%%%%%%%%
\section{Methodology}
\label{sec:method}

% TODO: Describe the skill-based framework design.
%   3.1 Label taxonomy and ontology
%   3.2 Skill file structure (TEAM.md, per-skill SKILL.md)
%   3.3 Triage, sprint-plan, rebalance, metrics workflow
%   3.4 Prototype application (FastAPI + React + GitHub API)
%   3.5 Evaluation approach (artifact quality, routing correctness, overhead)

%%%%%%%%%%%%%%%%%%%%%%%%%%%%%%%%%%%%%%%%%%%%%%%%%%%%%%%%%%%%%%%%%%%%%%%%%%%%%%%
\section{Findings and Analysis}
\label{sec:findings}

% TODO: Present results from prototype validation.
%   4.1 Artifact completeness and consistency
%   4.2 Routing and triage accuracy
%   4.3 Sprint planning quality
%   4.4 Human oversight touchpoints and governance
%   4.5 Limitations observed

%%%%%%%%%%%%%%%%%%%%%%%%%%%%%%%%%%%%%%%%%%%%%%%%%%%%%%%%%%%%%%%%%%%%%%%%%%%%%%%
\section{Conclusion}
\label{sec:conclusion}

% TODO: Summarize contributions. Reflect on practical applicability.
% Future work: multi-repo support, automated evaluation benchmarks,
% integration with CI/CD signals, fine-tuned triage models.

%%%%%%%%%%%%%%%%%%%%%%%%%%%%%%%%%%%%%%%%%%%%%%%%%%%%%%%%%%%%%%%%%%%%%%%%%%%%%%%
\bibliographystyle{IEEEtran}
% \bibliography{references}   % uncomment when references.bib is added

\end{document}
